% =====================================================================
%  ESTILO  -  Libro de Estilo TFG/TFM, EPS Universidad de Alicante
% =====================================================================

% --- Idioma y codificación -------------------------------------------
\usepackage[T1]{fontenc}
\usepackage[utf8]{inputenc}
\usepackage[spanish,es-tabla,es-nodecimaldot]{babel} % "Tabla" en vez de "Cuadro"
\addto\captionsspanish{%
  \renewcommand{\contentsname}{Índice de contenidos}%
  \renewcommand{\listfigurename}{Índice de figuras}%
  \renewcommand{\listtablename}{Índice de tablas}%
  \renewcommand{\appendixname}{Anexo}%
}

% --- Tipografía (apdo. 1.2) ------------------------------------------
% Familia 2 (texto) y fórmulas: Palatino (matemáticas en Palatino itálica)
\usepackage[sc]{mathpazo}
\linespread{1.05}
% Familia 1 (títulos y bibliografía): Helvética
\usepackage[scaled=0.92]{helvet}
% Familia 3 (código, ancho fijo): Courier
\usepackage{courier}
\usepackage{microtype}

% --- Página, márgenes e interlineado (apdo. 1.3 y 1.4) ---------------
% Interior 3 cm, exterior 2,5 cm, superior/inferior 2,5 cm,
% cabecera y pie DENTRO del margen (includeheadfoot)
\usepackage[a4paper,
            inner=3cm, outer=2.5cm, top=2.5cm, bottom=2.5cm,
            includeheadfoot, headheight=15pt]{geometry}
\usepackage{setspace}
\onehalfspacing                          % interlineado 1,5
\usepackage{emptypage}                   % páginas en blanco sin cabecera

% --- Colores, gráficos ------------------------------------------------
\usepackage{xcolor}
\usepackage{graphicx}
\graphicspath{{imagenes/}}
\usepackage{tikz}
\usetikzlibrary{calc,babel}
\usepackage{booktabs}
\usepackage{array}
\usepackage{tabularx}
\newcolumntype{Y}{>{\raggedright\arraybackslash}X}
\usepackage{enumitem}
\usepackage{csquotes}

% --- Títulos de secciones en Familia 1 (Helvética) --------------------
\usepackage{titlesec}
\titleformat{\chapter}[display]
  {\sffamily\bfseries\huge}
  {\Large\chaptertitlename\ \thechapter}{8pt}{\Huge}
\titleformat{name=\chapter,numberless}[display]
  {\sffamily\bfseries\Huge}{}{0pt}{}
\titlespacing*{\chapter}{0pt}{0pt}{30pt}
\titleformat*{\section}{\sffamily\bfseries\Large}
\titleformat*{\subsection}{\sffamily\bfseries\large}
\titleformat*{\subsubsection}{\sffamily\bfseries\normalsize}
\setcounter{secnumdepth}{3}              % numeración arábiga 1.1.1
\setcounter{tocdepth}{2}

% Índices también con títulos en Familia 1
\usepackage{tocloft}
\renewcommand{\cftchapfont}{\sffamily\bfseries}
\renewcommand{\cftchappagefont}{\sffamily\bfseries}
\renewcommand{\cfttoctitlefont}{\sffamily\bfseries\Huge}
\renewcommand{\cftloftitlefont}{\sffamily\bfseries\Huge}
\renewcommand{\cftlottitlefont}{\sffamily\bfseries\Huge}

% --- Cabeceras diferenciadas páginas pares/impares (apdo. 1.4) --------
\usepackage{fancyhdr}
\pagestyle{fancy}
\fancyhf{}
\renewcommand{\chaptermark}[1]{\markboth{#1}{}}
\renewcommand{\sectionmark}[1]{\markright{\thesection\ #1}}
\fancyhead[LE]{\sffamily\small\nouppercase{\leftmark}}   % par: capítulo
\fancyhead[RO]{\sffamily\small\nouppercase{\rightmark}}  % impar: sección
\fancyfoot[LE,RO]{\sffamily\small\thepage}               % nº en lado exterior
\renewcommand{\headrulewidth}{0.4pt}
\fancypagestyle{plain}{\fancyhf{}%
  \fancyfoot[LE,RO]{\sffamily\small\thepage}%
  \renewcommand{\headrulewidth}{0pt}}

% --- Figuras y tablas (apdo. 4) ---------------------------------------
\usepackage[font=small,labelfont=bf,labelsep=colon]{caption}
% Nota obligatoria con la fuente bajo cada figura/tabla (apdo. 4.3)
\newcommand{\fuente}[1]{%
  \par\vspace{2pt}{\centering\footnotesize\textit{Fuente:} #1\par}}

% --- Código (Familia 3) ----------------------------------------------
\usepackage{listings}
\renewcommand{\lstlistingname}{Código}
\renewcommand{\lstlistlistingname}{Índice de códigos}
\lstset{
  basicstyle=\ttfamily\small\singlespacing,
  keywordstyle=\bfseries\color{blue!60!black},
  commentstyle=\itshape\color{gray},
  stringstyle=\color{red!50!black},
  numbers=left, numberstyle=\tiny\color{gray},
  frame=single, rulecolor=\color{gray!50},
  breaklines=true, showstringspaces=false,
  tabsize=2, captionpos=b,
  literate={á}{{\'a}}1 {é}{{\'e}}1 {í}{{\'i}}1 {ó}{{\'o}}1 {ú}{{\'u}}1
           {Á}{{\'A}}1 {É}{{\'E}}1 {Í}{{\'I}}1 {Ó}{{\'O}}1 {Ú}{{\'U}}1
           {ñ}{{\~n}}1 {Ñ}{{\~N}}1 {ü}{{\"u}}1 {¿}{{?`}}1 {¡}{{!`}}1
}

% --- Abreviaturas y siglas (apdo. 2.5) --------------------------------
\usepackage[printonlyused]{acronym}

% --- Bibliografía (apdo. 2.8): APA, en Familia 1 ----------------------
% Para sistema numérico (ISO 690): style=iso-numeric (paquete biblatex-iso690)
\usepackage[backend=biber, style=apa]{biblatex}
\addbibresource{bibliografia.bib}
\renewcommand*{\bibfont}{\sffamily\small}

% --- Enlaces (siempre lo último) --------------------------------------
\usepackage[colorlinks=true, linkcolor=black, urlcolor=blue!60!black, citecolor=blue!60!black]{hyperref}

% --- Páginas preliminares (apdo. 2.3 y 2.4) ---------------------------
% Una sola hoja, alineación derecha, distribuida verticalmente
\newenvironment{paginaderecha}
  {\cleardoublepage\thispagestyle{empty}\vspace*{\fill}\begin{flushright}\itshape}
  {\end{flushright}\vspace*{\fill}\clearpage}

% Capítulo sin numerar que sí aparece en el índice
\newcommand{\capitulosinnumero}[1]{%
  \chapter*{#1}\addcontentsline{toc}{chapter}{#1}\markboth{#1}{#1}}

% añadido
\usepackage{float}